\section{Introduction}\label{sec:intro}

Every evening, an e-commerce fulfilment organisation commits money it cannot easily take back.
Line-haul trucks are tendered to carriers for the next night.
Labour shifts are scheduled at every fulfilment centre.
These commitments are made against a forecast; the orders they will serve have not been placed yet.
The next day, millions of units arrive one by one and each must be assigned, irrevocably, to an FC that has the item, the labour and space on a truck, and ideally a lane short enough to keep the next-day promise.
When the committed capacity is in the wrong place, the network pays for spot trucks, overtime, late deliveries and broken promises.

This report builds the complete chain of models behind these decisions for a realistically dimensioned network covering mainland Spain and Portugal:
65 demand zones, 10 FCs, 12 sort centres, 4 product categories and about 450,000 units per day, rising beyond a million on Black Friday and Prime Day.
The aim is to treat the three layers, \emph{forecast} $\rightarrow$ \emph{plan} $\rightarrow$ \emph{assign}, as one system and to evaluate each layer by what it costs the next one.

\begin{keybox}
\textbf{Executive summary.}
\begin{itemize}[leftmargin=1.2em]
  \item \textbf{Forecasting.} A global LightGBM quantile model with event-aligned features halves the day-ahead error of a seasonal-naive forecast (wMAPE 24.3\% $\to$ 11.9\%), and beats a zero-shot foundation model (Chronos-Bolt, 15.4\%). Rolling conformal calibration fixes its overconfidence (80\% intervals: 71.6\% $\to$ 79.3\% coverage). A Gaussian copula on forecast residuals makes the \emph{national} uncertainty honest: the coverage of a 90\% interval rises from 30\% to 82\%.
  \item \textbf{Planning.} The day-ahead truck and labour plan is a two-stage stochastic MIP. It has 130 integer variables and a recourse LP per scenario, i.e.\ 254{,}000 columns at 100 scenarios. Benders decomposition with an LP phase, in-out stabilisation and a partial-decomposition (Jensen) block solves it to 0.05\% in seconds: at least $90\times$ faster than classic Benders and $12\times$ faster than the extensive form at 100 scenarios. Sample average approximation (SAA) bounds certify the solutions to within 0.4\% of the true optimum (95\% confidence).
  \item \textbf{Value of forecasts.} Scored on the true demand distribution over 31 test days, the remaining forecast error costs \EUR{16}k per day (about \EUR{6}M per year). A per-series P90 buffer, a common heuristic, wastes \EUR{96}k per day. The stochastic plan beats the deterministic plan by \EUR{3}--\EUR{10}k per day when same-day recourse is cheap, and by \EUR{39}k per day when it is expensive. In that regime copula scenarios are worth \EUR{20}k per day more than independent ones, and the hedge flips from under- to over-provisioning exactly where newsvendor theory predicts.
  \item \textbf{Speed vs.\ cost vs.\ CO$_2$.} The Pareto frontier puts a price on the next-day promise. Raising the next-day share from 83.5\% to 91\% costs only \EUR{8}k per day, after which the frontier turns vertical at a ceiling set by the network. Speed is bought by routing, not by extra capacity, and it raises line-haul CO$_2$ by 8\%.
  \item \textbf{Online assignment.} Assigning orders greedily costs 7.1\% more than a clairvoyant LP. LP-dual bid prices re-solved hourly close 67\% of the gap, and fluid LP tracking closes 79\%: \EUR{110}k per day.
  \item \textbf{Investment what-if.} A fulfilment centre near Lisbon raises next-day delivery in northern Portugal by up to 85 percentage points, cuts CO$_2$ by 10\% and saves \EUR{9.3}M per year in operating cost (\EUR{18.3}M including the value of speed).
\end{itemize}
\end{keybox}

\paragraph{Contributions.}
\begin{enumerate}
  \item An end-to-end, reproducible forecast-to-fulfilment pipeline: a generative demand model with known ground truth, probabilistic forecasting with conformal calibration and copula scenarios, a two-stage stochastic capacity planner, and an online LP-based controller, all evaluated in one consistent cost model.
  \item A Benders implementation for the planning problem that combines warm-started parallel subproblems, an LP phase, in-out stabilisation, partial decomposition justified by a Jensen bound (\cref{prop:jensen}) and an adaptive master tolerance. An ablation study isolates the contribution of each component (\cref{sec:computation}).
  \item A \emph{decision-focused} evaluation of forecasts. Because the data-generating process is known, every plan is scored against the true conditional demand distribution. This prices forecast errors, buffer heuristics and dependence modelling in euros per day, and shows \emph{when} each matters.
  \item Multi-objective and strategic analyses with the same machinery: exact cost--speed and cost--CO$_2$ frontiers (including unsupported Pareto points that a weighted sum misses), and the business case of a new fulfilment centre.
\end{enumerate}

\paragraph{Outline.}
\begin{itemize}
  \item \Cref{sec:related} reviews related work.
  \item \Cref{sec:setting} describes the network, the costs and the demand model.
  \item \Cref{sec:forecasting} develops the forecasts and the scenario generator.
  \item \Cref{sec:planning} formulates the stochastic planning problem and the decomposition, which \cref{sec:computation} benchmarks computationally.
  \item \Cref{sec:planresults} reports the planning results: statistical quality, value of the stochastic solution, value of forecasts.
  \item \Cref{sec:pareto} covers the multi-objective analysis, \cref{sec:online} the online assignment, and \cref{sec:network} the network-design what-if.
  \item \Cref{sec:discussion} discusses limitations and next steps.
\end{itemize}
All code, data generation and experiments are available and reproducible from a single \texttt{make} target (\cref{app:repro}).
